\documentclass[../../main.tex]{subfiles}% 注意这里的文件路径不能用 ./main.tex ,否则用latexmk编译子文件会报错
\graphicspath{{\subfix{./image/}}{\subfix{../../image/}}} % 指定图片引用路径,后续可以直接使用图片文件名
% 如果图片存放在上级目录,则使用 ../../image/ 作为备用路径

% 例如:
% \begin{figure}[H]
% \centering
% \includegraphics[width=0.3\textwidth]{图.png}
% \caption{}
% \label{figure:图}
% \end{figure}
% 注意:上述\label{}一定要放在\caption{}之后,否则引用图片序号会只会显示??.

\begin{document}
% SOURCE p=19; printed=4; object=ch00-sec02-title; region=bottom
\Needspace{12\baselineskip}
\section{序关系}\label{section:ch00-orderings}
% SOURCE p=19,20; printed=4,5; object=ch00-partial-order; region=bottom,top

\begin{definition}[偏序]\label{definition:ch00-partial-order}
非空集合 $X$ 上的\textbf{偏序}是 $X$ 上具有下列性质的关系 $R$：
\begin{itemize}
\item 若 $xRy$ 且 $yRz$，则 $xRz$；
\item 若 $xRy$ 且 $yRx$，则 $x=y$；
\item 对所有 $x$ 都有 $xRx$。
\end{itemize}
\end{definition}
% SUPPLEMENT: id=v14-ch00-partial-order-remark; parent=ch00-partial-order; kind=motivation+hypotheses; origin=independent-derivation
\begin{remark}
\textbf{动机与条件说明（重构新增）。}偏序抽取“可以放在另一对象之前或之内”的比较规则，而允许有些对象无法比较。传递性使连续两次比较能合并；反对称性排除不同元素彼此都在对方之前；反身性使每个元素可与自身比较，适合非严格的序记号。三个条件彼此独立：在 $\{0,1,2\}$ 上，取所有对角有序对，再加 $(0,1)$、$(1,2)$，所得关系反身且反对称，却不传递；在 $\{0,1\}$ 上取全部有序对，关系反身且传递，却不反对称；非空集合上的空关系传递且反对称，却不反身。这些例子中未列出的有序对均不属于相应关系。本书规定偏序集非空，这是术语约定；三个关系公理本身也允许空集合，但空集合没有极大元，使用存在性结论时必须另行处理。
\end{remark}
% END SUPPLEMENT: id=v14-ch00-partial-order-remark


% SOURCE p=20; printed=5; object=ch00-linear-order; region=upper

\begin{definition}[线性序]\label{definition:ch00-linear-order}
若 $R$ 还满足
\begin{itemize}
\item 若 $x,y\in X$，则 $xRy$ 或 $yRx$，
\end{itemize}
就称 $R$ 为\textbf{线性序}（或\textbf{全序}）。
\end{definition}
% SUPPLEMENT: id=v14-ch00-linear-order-remark; parent=ch00-linear-order; kind=motivation+hypotheses; origin=independent-derivation
\begin{remark}
\textbf{动机与条件说明（重构新增）。}线性序在偏序的三条公理上增加任意两元素的可比较性，从而可以把所有元素安排在一条次序线上。“或”允许两个关系同时成立，但由反对称性可知这只发生在 $x=y$ 时。可比较性并非偏序自动具有：在 $\mathscr{P}(\{0,1\})$ 中，$\{0\}$ 与 $\{1\}$ 在包含关系下不可比较。这里没有要求存在首项、末项或相邻元素；这些都是另外的性质。
\end{remark}
% END SUPPLEMENT: id=v14-ch00-linear-order-remark


\begin{examples}\label{examples:ch00-order-examples}
例如，对任意集合 $E$，$\mathscr{P}(E)$ 按包含关系成为偏序集，而 $\mathbb{R}$ 按通常的序关系成为线性序集。
\end{examples}
仿照后一个例子，我们通常用 $\leqslant$ 表示偏序，并用 $x<y$ 表示 $x\leqslant y$ 但 $x\neq y$。注意，$X$ 上的偏序自然地在 $X$ 的每个非空子集上诱导出偏序。
% SUPPLEMENT: id=v14-ch00-order-examples-proof; parent=ch00-order-examples; kind=proof; origin=independent-derivation
\begin{proof}
\textbf{例子和限制断言的补写验证（重构新增）。}对 $\mathscr{P}(E)$ 上的包含关系，若 $A\subseteq B\subseteq C$，则每个 $x\in A$ 先属于 $B$，再属于 $C$，故 $A\subseteq C$；若同时 $A\subseteq B$ 与 $B\subseteq A$，则二者成员完全相同，故 $A=B$；每个 $A$ 显然满足 $A\subseteq A$。而 $\mathscr{P}(E)$ 总含空集，故为非空偏序集。实数通常的序关系由实数系的序公理满足这三条性质及任意两数可比较，故为线性序；这里以读者已知的实数序公理为基础。

若 $A\subseteq X$ 非空，限制关系为 $R\cap(A\times A)$。对 $x,y,z\in A$，原关系的传递性和反对称性仍然适用，而且传递所得的有序对仍在 $A\times A$ 中；每个 $x\in A$ 也仍满足 $xRx$。因此该限制确为偏序。如果原序是线性序，则任意 $x,y\in A$ 在 $X$ 中可比较，限制后仍可比较，故限制也是线性序。
\end{proof}
% END SUPPLEMENT: id=v14-ch00-order-examples-proof

\begin{definition}[序同构]\label{definition:ch00-order-isomorphism}
若两个偏序集 $X$ 和 $Y$ 之间存在双射 $f:X\to Y$，使得 $x_1\leqslant x_2$ 当且仅当 $f(x_1)\leqslant f(x_2)$，就称它们\textbf{序同构}。
\end{definition}
% SUPPLEMENT: id=v14-ch00-order-isomorphism-remark; parent=ch00-order-isomorphism; kind=motivation+hypotheses; origin=independent-derivation
\begin{remark}
\textbf{动机与条件说明（重构新增）。}序同构表示更换元素名称以后，比较关系完全不变。双射保证没有合并元素或遗漏目标元素；“当且仅当”中的正向保证保持已有比较，反向保证不会在目标中凭空增加原先没有的比较。仅有双射且保序仍不够：把二元素集合 $\{a,b\}$ 上只允许自身比较的偏序，恒等地映到增加 $a<b$ 的线性序，得到保序双射，却不反映比较，故不是序同构。事实上，“当且仅当”加上两端的反对称性已经蕴含单射：若 $f(x)=f(y)$，由目标反身性和反向蕴含得到 $x\leqslant y$、$y\leqslant x$，所以 $x=y$。显式写“双射”能直接表达完整对应，并强调满射仍是必需的；例如 $\mathbb{N}\hookrightarrow\mathbb{Z}$ 保持且反映通常的序，但遗漏非正整数。
\end{remark}
% END SUPPLEMENT: id=v14-ch00-order-isomorphism-remark



% SOURCE p=20; printed=5; object=ch00-maximal-bounds; region=middle

\begin{definition}[极大元与极小元]\label{definition:ch00-maximal-bounds}
若 $X$ 按 $\leqslant$ 成为偏序集，则 $X$ 的\textbf{极大元}（相应地，\textbf{极小元}）是这样的元素 $x\in X$：满足 $x\leqslant y$（相应地，$x\geqslant y$）的 $y\in X$ 只有 $x$ 本身。
\end{definition}
% SUPPLEMENT: id=v14-ch00-maximal-remark; parent=ch00-maximal-bounds; kind=motivation+hypotheses; origin=independent-derivation
\begin{remark}
\textbf{动机与条件说明（重构新增）。}极大元表示在当前偏序中无法再向严格更大的元素推进。条件 $x\in X$ 不可省略，否则谈的就不是集合自身的元素；对 $y\in X$ 的量化则确定了在哪个范围内检验极大性。条件只管与 $x$ 可比较且在其上的元素，并不要求每个元素都在 $x$ 之下。因此极大元与“比所有元素都大”的最大元含义不同；极小元与最小元也有同样区别。
\end{remark}
% END SUPPLEMENT: id=v14-ch00-maximal-remark


\begin{remark}
极大元与极小元可能存在，也可能不存在；除非序关系是线性序，否则它们未必唯一。
\end{remark}
% SUPPLEMENT: id=v14-ch00-maximal-examples-proof; parent=ch00-maximal-bounds; kind=proof; origin=independent-derivation
\begin{proof}
\textbf{存在性、唯一性断言的补写验证（重构新增）。}实数通常序下，任给 $x$，都有 $x-1<x<x+1$，所以既没有极大元，也没有极小元。单元素偏序集的唯一元素同时是极大元和极小元。在含两个不同元素的集合上只取相等关系，则每个元素既极大又极小，说明一般不能保证唯一。

在线性序集中，若 $x,y$ 都是极大元，则二者可比较。若 $x\leqslant y$，$x$ 的极大性给出 $y=x$；若 $y\leqslant x$，$y$ 的极大性也给出 $x=y$。对两个极小元，若 $x\leqslant y$，则 $y$ 的极小性给出 $x=y$；另一种次序由 $x$ 的极小性给出同样结论。因此两种元素在各自存在时均唯一。
\end{proof}
% END SUPPLEMENT: id=v14-ch00-maximal-examples-proof


\begin{definition}[上界与下界]\label{definition:ch00-upper-lower-bounds}
若 $E\subseteq X$，则 $E$ 的\textbf{上界}（相应地，\textbf{下界}）是这样的元素 $x\in X$：对所有 $y\in E$，都有 $y\leqslant x$（相应地，$x\leqslant y$）。
\end{definition}
% SUPPLEMENT: id=v14-ch00-bounds-remark; parent=ch00-upper-lower-bounds; kind=motivation+hypotheses; origin=independent-derivation
\begin{remark}
\textbf{动机与条件说明（重构新增）。}界把一整组元素同时控制在某一方向上。必须说明环境偏序集 $X$：上界或下界从 $X$ 中寻找，但不要求落在被控制的子集 $E$ 中。量词“对所有 $y\in E$”不能换为“某个”，否则不能控制整个集合。空子集也允许；此时全称条件没有反例，所以每个 $x\in X$ 同时是其上界和下界。若 $E$ 非空，能否找到界则取决于序关系和所选环境。
\end{remark}
% END SUPPLEMENT: id=v14-ch00-bounds-remark


\begin{remark}
$E$ 的上界未必是 $E$ 的元素；而且除非 $E$ 是线性序集，否则 $E$ 的极大元未必是 $E$ 的上界。（读者应当自己举一些例子。）
\end{remark}
% SUPPLEMENT: id=v14-ch00-bounds-examples-proof; parent=ch00-maximal-bounds; kind=proof; origin=independent-derivation
\begin{proof}
\textbf{原文留给读者的例子与验证（重构新增）。}在 $X=\mathbb{R}$ 中取 $E=(0,1)$，则 $1$ 是上界却不属于 $E$，$0$ 是下界也不属于 $E$。在 $X=\mathscr{P}(\{0,1\})$ 的包含偏序中，取 $E=\{\{0\},\{1\}\}$。两个成员都在 $E$ 中极大，因为另一个成员并不包含它；但它们都不是 $E$ 的上界，因为都不包含另一个成员。环境中的 $\{0,1\}$ 是 $E$ 的上界。

若 $E$ 的诱导序是线性的，而 $x\in E$ 是极大元，则任给 $y\in E$，或者 $y\leqslant x$，或者 $x\leqslant y$；后一种情形由极大性给出 $x=y$，仍有 $y\leqslant x$。故 $x$ 是 $E$ 的上界。这也验证了原文关于线性序的例外情形。
\end{proof}
% END SUPPLEMENT: id=v14-ch00-bounds-examples-proof



% SOURCE p=20; printed=5; object=ch00-well-order; region=middle

\begin{definition}[良序]\label{definition:ch00-well-order}
若 $X$ 按 $\leqslant$ 成为线性序集，并且 $X$ 的每个非空子集都有一个（必然唯一的）极小元，就称 $X$ 关于 $\leqslant$ \textbf{良序}，并且（不顾英语语法的规则）称 $\leqslant$ 为 $X$ 上的\textbf{良序}。
\end{definition}
% SUPPLEMENT: id=v14-ch00-well-order-remark; parent=ch00-well-order; kind=motivation+hypotheses; origin=independent-derivation
\begin{remark}
\textbf{动机与条件说明（重构新增）。}良序保证任何尚未处理的非空部分都有一个最先元素，因而能支持逐步选择和最小反例论证。必须要求“每个非空子集”，仅要求整个集合有最小元还不够：通常序下 $[0,\infty)$ 有最小元 $0$，但子集 $(0,1)$ 没有最小元，因为对任意 $x\in(0,1)$，$x/2$ 是其中更小的元素。非空限定排除空集，它没有任何元素可供选作极小元。线性序使极小元就是最小元且唯一：若 $m$ 极小，对任意 $y$ 的可比较性以及极小性给出 $m\leqslant y$；两个这样的元素再由反对称性相等。若删去线性序要求，二元素相等关系下每个非空子集都有极小元，但整个二元素集没有唯一的首项，无法表达这里的良序。
\end{remark}
% END SUPPLEMENT: id=v14-ch00-well-order-remark


\begin{example}\label{example:ch00-natural-well-order}
例如，$\mathbb{N}$ 按自然序成为良序集。
\end{example}
% SUPPLEMENT: id=v14-ch00-natural-well-order-proof; parent=ch00-natural-well-order; kind=proof; origin=independent-derivation
\begin{proof}
\textbf{补写验证（重构新增）。}通常的整数序使 $\mathbb{N}$ 成为线性序集。用自然数归纳法证明：$\{1,\ldots,n\}$ 的每个非空子集都有最小元。$n=1$ 时唯一的非空子集为 $\{1\}$。由 $n$ 到 $n+1$ 时，若子集与 $\{1,\ldots,n\}$ 相交，其交的最小元由归纳假设存在，并且不大于新增的 $n+1$；若不相交，则该子集只能是 $\{n+1\}$。归纳结论成立。现取非空 $A\subseteq\mathbb{N}$，取一个 $n\in A$。有限非空集 $A\cap\{1,\ldots,n\}$ 有最小元 $m$；$A$ 的其余元素都大于 $n\geqslant m$，故 $m$ 是 $A$ 的最小元。这证明 $\mathbb{N}$ 良序。
\end{proof}
% END SUPPLEMENT: id=v14-ch00-natural-well-order-proof



% SOURCE p=20; printed=5; object=ch00-maximal-principle-intro; region=middle-lower
现在陈述集合论的一条基本原理，并推出它的一些结论。

% SOURCE p=20; printed=5; object=ch00-hausdorff-maximal; number=0.1; region=lower
\begin{principle}[Hausdorff 极大原理]\label{principle:ch00-hausdorff-maximal}
每个偏序集都有一个极大的线性序子集。
\end{principle}
% SUPPLEMENT: id=v14-ch00-hausdorff-remark; parent=ch00-hausdorff-maximal; kind=motivation+hypotheses; origin=independent-derivation
\begin{remark}
\textbf{动机与条件说明（重构新增）。}线性序子集也称链；这里的“极大”按子集包含关系理解，意思是不能再添入元素而仍保持链，并非要求该链包含其他所有链。偏序提供链中比较的一致规则；不需要全体 $X$ 已经线性有序，也不要求链有上界。按本书约定 $X$ 非空，至少有单元素链可作为出发点。原理解决的是无限扩展过程可能没有最后一步的问题。它在本节被选作基本原则，其与选择公理的等价性将在本节补写证明中说明；不能把这种相对推导误称为不使用任何选择原则的绝对证明。
\end{remark}
% END SUPPLEMENT: id=v14-ch00-hausdorff-remark

% SOURCE p=20; printed=5; object=ch00-maximal-principle-explanation; region=lower

\begin{remark}
更具体地说，这意味着：若 $X$ 按 $\leqslant$ 成为偏序集，则存在集合 $E\subseteq X$，它关于 $\leqslant$ 是线性序的，且 $X$ 中任何真包含 $E$ 的子集都不是关于 $\leqslant$ 的线性序集。
\end{remark}
这条原理的另一个版本如下。

% SOURCE p=20; printed=5; object=ch00-zorn; number=0.2; region=lower
\begin{lemma}[Zorn 引理]\label{lemma:ch00-zorn}
若 $X$ 是偏序集，且 $X$ 的每个线性序子集都有上界，则 $X$ 有极大元。
\end{lemma}
% SUPPLEMENT: id=v14-ch00-zorn-remark; parent=ch00-zorn; kind=motivation+hypotheses; origin=independent-derivation
\begin{remark}
\textbf{动机与条件说明（重构新增）。}Zorn 引理把“每条逐步扩展的链都能被某个对象控制”转化为“存在不能再扩展的对象”。非空性来自偏序集的约定；若允许空偏序集且只检验非空链，假设会空真而结论失败。上界必须在 $X$ 内，否则无法把它当作候选极大元；不要求它已属于该链。要求所有链，而非只要求有限链：$\mathbb{N}$ 的每个有限子集有上界，但 $\mathbb{N}$ 自身没有上界，也没有极大元。链有上界是充分条件，不是存在极大元的必要条件。例如在 $\mathbb{N}\cup\{\ast\}$ 上保留自然数的通常序，只让 $\ast$ 与自身比较，则 $\ast$ 极大，而自然数链仍无上界。
\end{remark}
% END SUPPLEMENT: id=v14-ch00-zorn-remark

% SOURCE p=20; printed=5; object=ch00-hausdorff-zorn-equivalence; region=bottom

\begin{remark}
显然，\hyperref[principle:ch00-hausdorff-maximal]{Hausdorff 极大原理}蕴含\hyperref[lemma:ch00-zorn]{Zorn 引理}：$X$ 的一个极大线性序子集的上界就是 $X$ 的极大元。也不难看出，\hyperref[lemma:ch00-zorn]{Zorn 引理}蕴含\hyperref[principle:ch00-hausdorff-maximal]{Hausdorff 极大原理}。（将\hyperref[lemma:ch00-zorn]{Zorn 引理}应用于 $X$ 的线性序子集所组成的族，该族按包含关系成为偏序集。）
\end{remark}
% SUPPLEMENT: id=v14-ch00-hausdorff-zorn-proof; parent=ch00-hausdorff-zorn-equivalence; kind=proof; origin=independent-derivation
\begin{proof}
\textbf{两个方向的完整补写证明（重构新增）。}先假设 Hausdorff 极大原理，并设 $X$ 满足 Zorn 引理的假设。取极大链 $C\subseteq X$；因为 $X$ 非空，空链可以添入元素，所以 $C$ 非空。由假设，$C$ 有上界 $u\in X$。任给 $c\in C$ 都有 $c\leqslant u$，故 $C\cup\{u\}$ 仍是链，极大性给出 $u\in C$。若 $u\leqslant v$ 且 $v\in X$，则传递性使每个 $c\in C$ 都满足 $c\leqslant v$，所以 $C\cup\{v\}$ 也是链。再由极大性得 $v\in C$，而 $u$ 为上界给出 $v\leqslant u$。反对称性于是给出 $u=v$，证明 $u$ 极大。

反之，假设 Zorn 引理，设 $X$ 是非空偏序集。令 $\mathscr{C}$ 为 $X$ 的所有链组成的集合，允许加入空链；它是 $\mathscr{P}(X)$ 的子集，按包含关系为非空偏序集。设 $\mathscr{D}\subseteq\mathscr{C}$ 按包含关系成为链。若 $\mathscr{D}$ 非空，置 $U=\bigcup_{C\in\mathscr{D}}C$。任取 $x,y\in U$，存在 $C_x,C_y\in\mathscr{D}$ 分别含它们；两者按包含关系可比较，较大者同时含 $x,y$，所以 $x,y$ 在 $X$ 中可比较。故 $U$ 是链，并且包含 $\mathscr{D}$ 的所有成员，是它在 $\mathscr{C}$ 中的上界。若 $\mathscr{D}$ 为空，空链就是一个上界。因此 Zorn 引理适用于 $\mathscr{C}$，给出极大成员 $C_0$。它恰是 $X$ 中不能真扩大的链，得到 Hausdorff 极大原理。
\end{proof}
% END SUPPLEMENT: id=v14-ch00-hausdorff-zorn-proof



% SOURCE p=20; printed=5; object=ch00-well-ordering-principle; number=0.3; region=bottom
\begin{principle}[良序原理]\label{principle:ch00-well-ordering}
每个非空集合 $X$ 都可以良序化。
\end{principle}
% SUPPLEMENT: id=v14-ch00-well-ordering-remark; parent=ch00-well-ordering-principle; kind=motivation+hypotheses; origin=independent-derivation
\begin{remark}
\textbf{动机与条件说明（重构新增）。}良序原理允许给任意集合配上一套可逐项处理的次序，即使这个次序与已有结构无关。它断言存在某个良序，并不声称原有序就是良序；例如 $\mathbb{R}$ 的通常序没有最小元。$X$ 是集合这一前提保证可在其子集和关系的集合中进行构造。本书的非空约定排除了空良序；在允许空良序的常见扩展约定下，空集也平凡地满足结论。下方证明通过链的相容并使用 Zorn 引理，并不预先对每个子集任意指定一个良序。
\end{remark}
% END SUPPLEMENT: id=v14-ch00-well-ordering-remark

% SOURCE p=21; printed=6; object=ch00-proof-well-ordering; region=top
\begin{proof}
设 $\mathscr{W}$ 是 $X$ 的各个子集上的良序构成的族，在 $\mathscr{W}$ 上如下定义偏序。若 $\leqslant_1$ 和 $\leqslant_2$ 分别是子集 $E_1$ 和 $E_2$ 上的良序，则当下列条件成立时，规定在此偏序中 $\leqslant_1$ 先于 $\leqslant_2$：
\begin{enumerate}[label=(\roman*),ref=(\roman*)]
\item\label{item:ch00-well-order-extension} $\leqslant_2$ 延拓 $\leqslant_1$，即 $E_1\subseteq E_2$，且 $\leqslant_1$ 与 $\leqslant_2$ 在 $E_1$ 上一致；
\item\label{item:ch00-well-order-end-extension} 若 $x\in E_2\setminus E_1$，则对所有 $y\in E_1$ 都有 $y\leqslant_2 x$。
\end{enumerate}
读者可以验证，\hyperref[lemma:ch00-zorn]{Zorn 引理}的假设成立，因而 $\mathscr{W}$ 有一个极大元。这个极大元必定是 $X$ 本身上的良序，因为若 $\leqslant$ 是 $X$ 的真子集 $E$ 上的良序，且 $x_0\in X\setminus E$，则只需规定对所有 $x\in E$ 都有 $x\leqslant x_0$，便可将 $\leqslant$ 延拓为 $E\cup\{x_0\}$ 上的良序。
\end{proof}
% SUPPLEMENT: id=v14-ch00-well-ordering-proof; parent=ch00-proof-well-ordering; kind=proof+motivation+hypotheses; origin=independent-derivation
\begin{definition*}\phantomsection\label{definition:ch00-initial-segment-supplement}
\textbf{补充用语（重构新增）。}在本节补写中，线性序集 $L$ 的子集 $I$ 称为一个\textbf{初始段}，若 $x\in I$、$y\in L$ 且 $y<x$ 蕴含 $y\in I$。此处允许空初始段和 $L$ 本身。
\end{definition*}
\begin{remark}
\textbf{动机与条件说明（重构新增）。}初始段保留每个已纳入元素之前的全部元素，因而能表达“只在尾部延长”。条件中的 $y\in L$ 固定了检验范围，$y<x$ 则限制为更早的元素；不要求把更晚的元素一并纳入。在线性序中，这等价于 $I$ 中每个元素都严格小于 $L\setminus I$ 中每个元素：若外部元素反而更早，下闭条件就会把它纳入；反之，这种内外比较直接排除遗漏更早元素。原证明的\hyperref[item:ch00-well-order-end-extension]{尾部条件}正是在已有良序上追加尾部的要求；\hyperref[item:ch00-well-order-extension]{延拓条件}保证原来的比较不被改变。本节的辅助用语允许全体 $L$；后文\cref{definition:ch00-initial-segments}按原书另将 $I_x=\{y:y<x\}$ 称为元素 $x$ 的初始段，那里不包含全体的情形，两种约定须区分。本节证明不依赖后文结果。
\end{remark}
\begin{proof}
\textbf{原证明中验证步骤的补全（重构新增）。}把 $\mathscr{W}$ 的成员精确记为 $(E,R)$，其中 $\varnothing\ne E\subseteq X$，$R\subseteq E\times E$ 是良序。它是 $\mathscr{P}(X)\times\mathscr{P}(X\times X)$ 的子集，故为集合。因 $X$ 非空，任取 $x\in X$，单元素集 $\{x\}$ 的唯一良序给出一个成员。

记原文定义的关系为 $\preccurlyeq$。它反身；若两成员互相先于对方，则其定义域互相包含而相等，比较关系由延拓条件一致，故成员相同，得到反对称性。验证传递性：设 $(E_1,R_1)\preccurlyeq(E_2,R_2)\preccurlyeq(E_3,R_3)$。关系延拓显然传递。若 $x\in E_3\setminus E_1$、$y\in E_1$，当 $x\in E_2$ 时，用第一段延拓的尾部条件再用关系一致性，得 $yR_3x$；当 $x\notin E_2$ 时，用第二段延拓的尾部条件也得 $yR_3x$。故 $\preccurlyeq$ 为偏序。

设 $\mathscr{D}$ 是其中的非空链，令 $E=\bigcup_{(E_i,R_i)\in\mathscr{D}}E_i$、$R=\bigcup_{(E_i,R_i)\in\mathscr{D}}R_i$。链中两个成员互为尾部延拓之一，因而在共同定义域上完全一致。任取有限多个 $E$ 中的元素，它们所在的有限多个成员中有一个按延拓次序最大，包含这些元素；在该成员中检验关系，即可得到 $R$ 的反身性、反对称性、传递性和线性可比较性。因此 $R$ 是 $E$ 上的线性序，且每个 $R_i$ 是其限制。

每个 $E_i$ 是 $E$ 的初始段：若 $x\in E_i$、$y\in E$ 且 $y<_R x$，取链中同时包含 $x,y$ 且延拓 $E_i$ 的成员；尾部条件不允许新增的 $y$ 早于 $x$，所以 $y\in E_i$。现取非空 $S\subseteq E$，选 $s\in S$ 及包含它的 $E_i$。良序 $R_i$ 给出 $S\cap E_i$ 的最小元 $m$。在 $S\cap E_i$ 内它最小；在 $S\setminus E_i$ 中的每个元素都位于整个初始段 $E_i$ 之后，所以也大于 $m$。故 $m$ 为 $S$ 的最小元，$R$ 是良序。以上同时说明 $(E,R)$ 延拓每个 $(E_i,R_i)$ 且满足尾部条件，因此是链的上界。空链则可用先前构造的任一单元素良序作上界。

Zorn 引理现在确实适用。设 $(E,R)$ 是极大成员。若 $E\ne X$，取 $x_0\in X\setminus E$，把 $x_0$ 放在 $E$ 的全部元素之后，同时保留 $R$。新增序显然线性；其任意非空子集若与 $E$ 相交，就取该交在 $R$ 下的最小元，否则只能为 $\{x_0\}$，取 $x_0$ 为最小元。因此新增序仍是良序，且严格尾部延拓 $(E,R)$，与极大性矛盾。故 $E=X$。
\end{proof}
\begin{remark}
\textbf{尾部条件的作用（重构新增）。}只要求良序彼此延拓而不要求初始段性质，相容并可能失去良序性。令 $E_n=\{-1,\ldots,-n\}$，各自取通常整数序。每个 $E_n$ 有限，因而良序，且后一个序延拓前一个；并集却是负整数集，没有最小元，因为每个 $-n$ 前面还有 $-(n+1)$。原证明的尾部条件排除的正是不断插入更早元素的情形。
\end{remark}
% END SUPPLEMENT: id=v14-ch00-well-ordering-proof


% SOURCE p=21; printed=6; object=ch00-choice; number=0.4; region=upper-middle
\begin{axiom}[选择公理]\label{axiom:ch00-choice}
若 $\{X_\alpha\}_{\alpha\in A}$ 是由非空集合组成的非空族，则 $\prod\limits_{\alpha\in A}X_\alpha$ 非空。
\end{axiom}
% SUPPLEMENT: id=v14-ch00-choice-remark; parent=ch00-choice; kind=motivation+hypotheses; origin=independent-derivation
\begin{remark}
\textbf{动机与条件说明（重构新增）。}选择公理把各处分别存在可选元素，提升为同时存在一份完整选择函数。每个因子非空是必要的：一个空因子就会使相应坐标无值可选，从而乘积为空。指标族非空则是原书的表述约定；若指标集为空，乘积含唯一空函数，结论无需使用选择公理也成立。因子不必互不相交，也不需要有统一的选取规则。有限个非空因子的选择可由有限次存在量词推理和归纳完成；任意指标族才是这里额外选择原则所处理的范围。
\end{remark}
% END SUPPLEMENT: id=v14-ch00-choice-remark

% SOURCE p=21; printed=6; object=ch00-proof-choice; region=upper-middle
\begin{proof}
令 $X=\bigcup_{\alpha\in A}X_\alpha$。在 $X$ 上选取一个良序，并对 $\alpha\in A$，令 $f(\alpha)$ 为 $X_\alpha$ 的极小元。于是 $f\in\prod\limits_{\alpha\in A}X_\alpha$。
\end{proof}
% SUPPLEMENT: id=v14-ch00-choice-proof; parent=ch00-proof-choice; kind=proof; origin=independent-derivation
\begin{proof}
\textbf{良定义性步骤的补全（重构新增）。}因 $A$ 非空，取 $\alpha_0\in A$；$X_{\alpha_0}$ 非空，所以总并 $X$ 非空，良序原理适用。每个非空 $X_\alpha\subseteq X$ 都有唯一最小元 $m_\alpha$：存在性由良序给出，唯一性由反对称性给出。以有序对 $(\alpha,m_\alpha)$ 组成的关系定义 $f$，每个指标恰有一个输出，故它是 $A\to X$ 的函数；而 $m_\alpha\in X_\alpha$ 恰是乘积要求的逐项条件。这一步依靠已固定良序确定唯一对象，不是在证明中再次暗中使用任意选择。
\end{proof}
% END SUPPLEMENT: id=v14-ch00-choice-proof


% SOURCE p=21; printed=6; object=ch00-choice-representatives; number=0.5; region=middle
\begin{corollary}\label{corollary:ch00-choice-representatives}
若 $\{X_\alpha\}_{\alpha\in A}$ 是由非空集合组成的互不相交的族，则存在集合 $Y\subseteq\bigcup_{\alpha\in A}X_\alpha$，使得对每个 $\alpha\in A$，$Y\cap X_\alpha$ 恰含一个元素。
\end{corollary}
% SUPPLEMENT: id=v14-ch00-representatives-remark; parent=ch00-choice-representatives; kind=motivation+hypotheses; origin=independent-derivation
\begin{remark}
\textbf{动机与条件说明（重构新增）。}结论要从每块中取出一个代表，并让这些代表组成一个集合。非空性保证每块确有代表；互不相交保证从另一块取出的代表不会再次落入当前块，从而将“至少一个”提升为“恰好一个”。它是该普遍结论的充分结构条件；去掉它后并非所有族都能做到。例如三块 $\{1,2\}$、$\{2,3\}$、$\{1,3\}$ 不存在共同的代表集合：若 $a_i$ 为 $i$ 是否属于 $Y$ 的 $0$ 或 $1$ 指示数，则三个恰一条件要求 $a_1+a_2=a_2+a_3=a_1+a_3=1$，相加得到偶数 $2(a_1+a_2+a_3)=3$，矛盾。某些相交族仍可能满足结论，所以这不是对每个具体族的必要条件。$Y$ 限制在总并内只为排除无关元素。
\end{remark}
% END SUPPLEMENT: id=v14-ch00-representatives-remark

% SOURCE p=21; printed=6; object=ch00-proof-choice-representatives; region=middle
\begin{proof}
取 $Y=f(A)$，其中 $f\in\prod\limits_{\alpha\in A}X_\alpha$。
\end{proof}
% SUPPLEMENT: id=v14-ch00-representatives-proof; parent=ch00-proof-choice-representatives; kind=proof; origin=independent-derivation
\begin{proof}
\textbf{遗漏细节及空族情形的补全（重构新增）。}若 $A=\varnothing$，取 $Y=\varnothing$ 即满足全部条件。若 $A\ne\varnothing$，由选择公理取原文中的 $f$，置 $Y=f(A)$。每个 $f(\alpha)\in X_\alpha$，故 $Y$ 包含于总并，且 $f(\alpha)\in Y\cap X_\alpha$。反之，若 $y\in Y\cap X_\alpha$，由 $y\in f(A)$ 存在 $\beta\in A$ 使 $y=f(\beta)\in X_\beta$；于是 $y\in X_\alpha\cap X_\beta$。互不相交条件迫使 $\beta=\alpha$，从而 $y=f(\alpha)$。所以 $Y\cap X_\alpha=\{f(\alpha)\}$，恰含一个元素。
\end{proof}
% END SUPPLEMENT: id=v14-ch00-representatives-proof

% SOURCE p=21; printed=6; object=ch00-choice-equivalence; region=middle
我们从\hyperref[principle:ch00-hausdorff-maximal]{Hausdorff 极大原理}推出了\hyperref[axiom:ch00-choice]{选择公理}；事实上，可以证明二者在逻辑上等价。
% SUPPLEMENT: id=v14-ch00-choice-implies-maximal-proof; parent=ch00-choice-equivalence; kind=proof; origin=independent-derivation
\begin{proof}
\textbf{选择公理推出极大原理的补写证明（重构新增）。}正向推导已在本节给出。为证明逆向，此处只假设选择公理，先推出 Zorn 引理，再使用上面已证明的 Zorn 引理蕴含 Hausdorff 极大原理的方向。以下构造不引用尚待推出的极大原理或良序原理。

设非空偏序集 $X$ 的每条链都有上界。反设 $X$ 没有极大元。对每条链 $C\subseteq X$，允许 $C=\varnothing$，令
\begin{align*}
U_C=\{x\in X:c<x\text{ 对每个 }c\in C\}.
\end{align*}
这些集合均非空：非空链有上界 $u$；由于 $u$ 不是极大元，存在 $v>u$，于是对每个 $c\in C$ 有 $c\leqslant u<v$，从传递性和反对称性得到 $c<v$。空链时 $U_C=X$。所有链组成 $\mathscr{P}(X)$ 的子集，因此 $\{U_C\}$ 是一个以集合为指标的非空集合族。由选择公理，固定函数 $h$，使 $h(C)\in U_C$ 对每条链成立。

考虑 $X$ 中所有满足下述条件的子集 $A$：$A$ 在 $X$ 的诱导序下良序，并且对每个 $a\in A$ 都有
\begin{align*}
a=h(A_{<a}),\qquad A_{<a}=\{x\in A:x<a\}.
\end{align*}
为使构造包括起点，在此也允许空子集，空集的条件按空真解释；这只是辅助构造的约定。把这些子集组成的集合记为 $\mathscr{T}$。它是 $\mathscr{P}(X)$ 的子集，且非空；例如 $\{h(\varnothing)\}$ 属于其中。下面证明任意两个成员互为初始段之一。

取 $A,B\in\mathscr{T}$。令 $D$ 为同时是 $A$ 与 $B$ 的初始段的所有集合的并。空集属于这类；这些集合的并仍是初始段，因为若 $x$ 属于其中一段而 $y<x$，则 $y$ 也属于该段。因此 $D$ 是两者的共同初始段，并且包含每个共同初始段。两边的序都是 $X$ 的限制，不存在关系不一致的问题。若 $D\ne A$ 且 $D\ne B$，良序性给出 $a=\min(A\setminus D)$ 和 $b=\min(B\setminus D)$。由于 $D$ 初始且 $a,b$ 各为第一个剩余元素，有 $A_{<a}=D=B_{<b}$；于是相容条件给出 $a=h(D)=b$。但这样 $D\cup\{a\}$ 又成为两者的共同初始段，严格包含 $D$，与其构造矛盾。因此 $D=A$ 或 $D=B$，即所需的初始段比较成立。

置 $T=\bigcup_{A\in\mathscr{T}}A$。任意两元素可放在同一个较大成员中，故 $T$ 是链。每个 $A\in\mathscr{T}$ 都是 $T$ 的初始段：若 $x\in A$、$y\in T$ 且 $y<x$，选含 $y$ 的成员 $B$；若 $B\subseteq A$ 则直接有 $y\in A$，否则 $A$ 是 $B$ 的初始段，也给出 $y\in A$。任取非空 $S\subseteq T$，选择 $s\in S$ 及含 $s$ 的成员 $A$。$S\cap A$ 的最小元 $m$ 存在；$S$ 在 $A$ 外的所有元素都在整个初始段 $A$ 之后，所以 $m$ 也是 $S$ 的最小元。故 $T$ 良序。对每个 $t\in T$，取含它的 $A$，由初始段性质得 $T_{<t}=A_{<t}$，所以 $t=h(T_{<t})$。因此 $T$ 本身也属于 $\mathscr{T}$。

现在 $z=h(T)$ 严格大于 $T$ 的全部元素，故 $z\notin T$。把 $z$ 加在 $T$ 的末端，得到的 $T'=T\cup\{z\}$ 仍良序：非空子集若交于 $T$，沿用该交的最小元；否则只能是 $\{z\}$。旧元素的前驱集合不变，新元素的前驱集合为 $T$，所以 $T'$ 的每个元素都满足规定的选择等式。于是 $T'\in\mathscr{T}$，按总并定义应有 $T'\subseteq T$，却又有 $z\in T'\setminus T$，矛盾。反设不成立，$X$ 有极大元，Zorn 引理得证。结合已经证明的两个蕴含方向，选择公理与 Hausdorff 极大原理等价。
\end{proof}
\begin{remark}
\textbf{补写证明的逻辑依赖（重构新增）。}逆向证明只在选取 $h(C)\in U_C$ 时使用选择公理；随后所有对象都由明确的集合构造给出。对良序子集所用的是它们各自的定义性质，没有先给整个 $X$ 配置良序，也没有再次引用 Zorn 引理寻找极大成员。因此这条逆向推导不会与本节先前从极大原理推出选择公理的推导形成循环。
\end{remark}
% END SUPPLEMENT: id=v14-ch00-choice-implies-maximal-proof

\end{document}
